\documentclass[12pt,a4paper,oneside]{article}
\usepackage[brazil]{babel}
\usepackage[utf8]{inputenc}
\usepackage{olymp}

\setcounter{problem}{1}

\begin{document}

\begin{problem}{Senha segura}{senha}{2}
 
Em certa empresa, uma senha é considerada segura se tiver letras maiúsculas, minúsculas, dígitos e um caractere especial (algum que não é letra nem dígito).

Por exemplo, a senha \texttt{a*D45x} seria segura, mas \texttt{R2d2} não, pois não contém caractere especial, nem \texttt{ufv\#23}, pois não contém letra maiúscula.

Se a senha tiver caracteres de apenas um tipo, ela é considerada ``fraca''. Se tiver caracteres de mais de um tipo, mas não de todos, é considerada ``media''.

%\Notes
%
%Observações, se tiver.

\InputFile

A entrada contém apenas uma linha, contendo uma senha. A senha pode conter letras (maiúsculas ou minúsculas), dígitos (de `0' a `9') e outros caracteres, mas não contém espaços. Restrições: a senha tem no máximo 30 caracteres.

\OutputFile

Seu programa deve gerar uma linha na saída, contendo uma das seguintes palavras, dependendo da segurança da senha: ``\texttt{fraca}'', ``\texttt{media}'' ou ``\texttt{segura}''
% \Notes
% Preste atenção aos tipos das variáveis, pois $F(90) \approx 2^{61}$.

\Example

\begin{example}
\exmp{
123456
}{
fraca
}%
\end{example}

\begin{example}
\exmp{
a*D45x
}{
segura
}%
\end{example}

\begin{example}
\exmp{
R2d2
}{
media
}%
\end{example}

\begin{example}
\exmp{
C-3PO
}{
media}%
\end{example}

\begin{example}
\exmp{
C-3po
}{
segura}%
\end{example}

%Comentários sobre o exemplo, se necessário.

\end{problem}

\end{document}
